\section{La reserva es un bien público: el cuarto precio}
\label{sec:standby}

\Cref{co:terminacion} garantiza terminación y no eficiencia, y
\cref{th:descomposicion} contabiliza la brecha sin acotarla. Esta sección
exhibe una fuente de ineficiencia \emph{estructural} —no un artefacto del
protocolo ni de la topología— y el mecanismo que la corrige. Es, además, un
resultado contra la propia arquitectura: la reserva que el sistema necesita no
emerge del juego tal como está planteado.

\subsection{Dos fallos distintos: el de incentivo y el de coordinación}

Antes de la reserva conviene separar dos formas de fracaso que se confunden, y
que exigen remedios distintos. La literatura de decisión estratégica insiste en
la diferencia entre el dilema del prisionero —donde desertar es dominante y
cooperar nunca es equilibrio— y la \emph{caza del ciervo}, donde el éxito exige
concurrencia y la elección de cada uno depende de lo que espere de los demás:
lanzarse a la presa mayor es arriesgado, porque quien coopera y no encuentra
compañía se queda con las manos vacías \cite{weiss2023games}.

\begin{proposicion}[El reclutamiento es una caza del ciervo]
\label{pr:caza-ciervo}
Sea un plantel de $n$ AMR, cada uno con la opción de unirse a una coalición de
tamaño requerido $r^\star$ o dedicarse a otra tarea con pago $w$. Unirse cuesta
$c$ y rinde $B$ si la coalición alcanza $r^\star$, y nada si no. Si
$B-c\ge w$ y $r^\star\ge2$, entonces \emph{tanto} el perfil en que todos se unen
\emph{como} aquel en que ninguno lo hace son equilibrios de Nash.
\end{proposicion}

\begin{proof}
Con todos dentro, desviarse da $w\le B-c$. Con nadie dentro, unirse en solitario
no alcanza $r^\star\ge2$ y da $-c<w$.
\end{proof}

\begin{observacion}[La cuenca de la cooperación se estrecha con el tamaño]
\label{obs:cuenca-ciervo}
La coexistencia de equilibrios no es lo relevante; lo es cuál se alcanza. Sea
$p^\star$ la creencia mínima sobre la participación ajena que hace de unirse la
mejor respuesta. Con $n=12$, $B=10$, $c=1$ y $w=3$ se obtiene
\[
  \begin{array}{c|ccccccccc}
    r^\star & 2 & 3 & 4 & 5 & 6 & 7 & 8 & 9 & 10\\\hline
    p^\star & 0{,}045 & 0{,}123 & 0{,}205 & 0{,}290 & 0{,}376 & 0{,}463 &
    0{,}551 & 0{,}641 & 0{,}732
  \end{array}
\]
La cuenca de la cooperación, $1-p^\star$, cae de $0{,}955$ a $0{,}268$. Esto es
un segundo motivo, independiente del físico, por el que las coaliciones grandes
son peores: \cref{co:agarre} da la razón mecánica —la movilidad se interseca— y
esta da la estratégica —cuanta más gente hace falta, más difícil es creer que
acudirá—. Ambas empujan en la misma dirección y se suman.
\end{observacion}

\figCiervo

\begin{observacion}[Por qué un precio no basta, y qué sí lo resuelve]
\label{obs:precio-no-basta}
Es tentador suponer que el instrumento del resto de este trabajo —un precio—
resuelve también esto. Solo en parte. Subvencionando la participación con $s$ y
con $r^\star=6$, la creencia crítica baja de $0{,}376$ a $0{,}338$ con $s=1$, a
$0{,}241$ con $s=3$, y solo con $s=4$ llega a cero, es decir, solo entonces
cooperar pasa a ser dominante y el mal equilibrio desaparece. Por debajo de ese
umbral la subvención \emph{estrecha} la cuenca del fracaso pero no la elimina.

La distinción importa porque el fallo no es de incentivo sino de expectativa: un
AMR que preferiría unirse se abstiene porque no confía en que los demás acudan, y
ningún pago moderado corrige una creencia. Lo que sí la corrige es el
compromiso verificable, y aquí conviene notar algo sobre la propia arquitectura:
la guarda de \textsf{REC} a \textsf{FORM} en \cref{def:automata} exige que la
coalición esté \emph{certificada} antes de avanzar, es decir, que los miembros
necesarios estén comprometidos. Esa guarda se introdujo por factibilidad física.
Cumple además una segunda función que no se había nombrado: convierte una
decisión simultánea bajo incertidumbre en una secuencial con compromiso
observable, que es precisamente el remedio de una caza del ciervo. El certificado
no solo comprueba que la coalición \emph{puede}; también hace creíble que
\emph{acudirá}.
\end{observacion}

\subsection{Por qué la teoría cooperativa estándar no se aplica sin más}

Antes de importar herramientas conviene comprobar la hipótesis que casi todas
comparten. Un juego es \emph{superaditivo} si unir dos grupos disjuntos vale al
menos lo que valen por separado, y esa condición es la que justifica, por
ejemplo, que el valor de Shapley prediga la formación de la gran coalición
\cite{okada2026games}.

\begin{observacion}[El juego de este trabajo no es superaditivo]
\label{obs:no-superaditivo}
Tomando como valor de una coalición la aceleración que puede imprimir a la
carga, corregida por la fracción de maniobras que sus contactos permiten sin
saturar la demanda angular —es decir, con la física de
\cref{obs:rigidez-alcance} y no con el bloqueo por rango que aquella observación
descarta—, el valor medio por tamaño resulta $0{,}430$, $0{,}486$, $0{,}293$,
$0{,}559$, $0{,}810$ y $1{,}053$ para coaliciones de uno a seis miembros: ni
monótono ni cóncavo. Sobre quinientos pares disjuntos, el $78\%$ viola la
superaditividad.

La razón es la ya conocida: la capacidad se suma pero la movilidad se interseca
(\cref{co:agarre}), y a eso se añade que cada miembro aporta masa. Unir dos
grupos bien orientados puede producir un conjunto mal orientado que valga menos
que cualquiera de los dos.

La consecuencia es de método: los resultados de teoría cooperativa que presuponen
superaditividad —entre ellos la predicción de la gran coalición— no son
aplicables aquí, y la formación de coaliciones \emph{parciales} no es un defecto
del mecanismo sino la respuesta correcta a un juego que no premia agrupar.
\end{observacion}

\subsection{La pregunta que la teoría cooperativa hace y este trabajo no hacía}

Todo lo anterior compara \emph{criterios} de reparto. La teoría de juegos
cooperativos plantea una pregunta anterior y más severa: dado un reparto,
¿existe algún subgrupo que, rompiendo y actuando por su cuenta, mejoraría a
\emph{todos} sus miembros? Si existe, se dice que ese subgrupo \emph{bloquea} el
reparto. El \emph{núcleo} es el conjunto de repartos que ningún subgrupo bloquea
\cite{okada2026games}.

\begin{proposicion}[El núcleo de este juego es vacío salvo en un caso]
\label{pr:nucleo}
Modélese la coalición como juego de umbral: $v(S)=B$ si $\lvert S\rvert\ge
r^\star$ y $v(S)=0$ en otro caso, que es la situación de una carga que no se
mueve sin un mínimo de miembros. Entonces el núcleo es no vacío si y solo si
$r^\star=n$, es decir, solo si hacen falta \emph{todos} los AMR disponibles.
\end{proposicion}

\begin{observacion}[Comprobado sobre el rango de operación]
\label{obs:nucleo-vacio}
Resolviendo el programa lineal que caracteriza el núcleo para todas las
coaliciones: con $n=4$ y $r^\star=4$ el núcleo es no vacío; con $r^\star=3$ y
$r^\star=2$, vacío. Con $n=6$: no vacío en $r^\star=6$, vacío en $4$ y en $2$.
Con $n=8$ y $r^\star=5$, vacío.

La lectura es incómoda y conviene no suavizarla: \emph{en el juego de umbral, en
cuanto sobra un AMR, ninguna repartición del valor resiste la ruptura de un
subgrupo}. No es un defecto del mecanismo de precios —ningún criterio lo
evitaría, porque el núcleo está vacío—; es una propiedad de ese juego. Y la
condición que lo salvaría, $r^\star=n$, es precisamente la que \cref{co:agarre} y
\cref{obs:cuenca-ciervo} desaconsejan por otras razones.
\end{observacion}

\begin{proposicion}[Donde el reparto \emph{sí} es estable]
\label{pr:nucleo-divisible}
Considérese en cambio el reparto de esfuerzo como modelo de producción
\emph{divisible}, $v(S)=\max\{c^{\!\top}y:\ Ay\le\sum_{i\in S}b_i,\ y\ge0\}$, con
óptimo finito y dualidad fuerte para la gran coalición. Sea $\lambda^\star\ge0$
una solución dual y defínase el pago $x_i=\lambda^{\star\top}b_i$. Entonces $x$
pertenece al núcleo.
\end{proposicion}

\begin{proof}
La dualidad fuerte da $\sum_i x_i=v(N)$. El mismo $\lambda^\star$ es dualmente
factible para cualquier subconjunto, porque $A$ y $c$ no cambian; por dualidad
débil, $v(S)\le\lambda^{\star\top}\sum_{i\in S}b_i=\sum_{i\in S}x_i$, de modo que
ninguna coalición bloquea.
\end{proof}

\begin{observacion}[Dos regímenes, no uno]
\label{obs:nucleo-dos-regimenes}
La comprobación numérica separa los dos casos sin ambigüedad. Sobre un modelo
divisible con cinco AMR, tres recursos y cuatro actividades, el reparto dual suma
exactamente $v(N)=32{,}691$ y \emph{ninguna} de las coaliciones propias lo
bloquea: el peor exceso $v(S)-\sum_{i\in S}x_i$ es $3{,}6\times10^{-15}$. Sobre
el juego de umbral con los mismos cinco AMR, el núcleo es vacío.

El alcance de \cref{obs:nucleo-vacio} es, por tanto, el juego de umbral. La arquitectura de
precios de este trabajo \emph{sí} produce un reparto estable frente a rupturas en
la capa donde el esfuerzo es divisible —que es el reparto de \emph{wrench}, el
apartado (a) de \cref{obs:alcance-poblacional}—, y no puede producirlo en la capa
de reclutamiento atómico, apartado (c). La vacuidad del núcleo es una propiedad
de la indivisibilidad, no del mecanismo.

Esto reordena también \cref{obs:nucleo-acoplamiento}: el acoplamiento rígido no
sustituye a un núcleo vacío en general, sino específicamente en la decisión
discreta de pertenencia. Una vez dentro, el precio dual basta.
\end{observacion}

\begin{observacion}[Y una salvedad sobre el propio modelo de juego]
\label{obs:funcion-particion}
Ambos cálculos suponen que el valor de una coalición depende solo de \emph{quiénes}
la forman, es decir, que el juego admite función característica. Con corredores
compartidos eso deja de ser cierto: si una coalición cambia de ruta modifica los
tiempos de las demás, de modo que el valor debe escribirse $v(S,\Pcal,\rho)$ con
la partición y el estado de recursos. Es un juego en \emph{función de partición},
y en esa clase ni la vacuidad del núcleo ni \cref{pr:nucleo-divisible} se
trasladan automáticamente. Sumar valores independientes de cargas que comparten
pasillo no está autorizado, y este trabajo no lo hace fuera de esta subsección.
\end{observacion}

\begin{observacion}[Lo que sostiene la coalición no es el reparto]
\label{obs:nucleo-acoplamiento}
De \cref{pr:nucleo} se sigue que la estabilidad de una coalición no puede
descansar en cómo se reparte el valor. Debe descansar en que la ruptura sea
\emph{inviable}, no en que sea poco atractiva. Y en este sistema lo es por
razones físicas: en \textsf{TRANS} los miembros comparten un cuerpo rígido, y
abandonarlo no es una decisión de cartera sino una maniobra que detiene la carga
y activa las barreras de \cref{th:astop-cuerpo}.

Es decir, el acoplamiento rígido no cumple solo una función mecánica: es el
dispositivo que sustituye a un núcleo vacío en la decisión de pertenencia
(véase la matización de \cref{obs:nucleo-dos-regimenes}). Donde la teoría cooperativa exige
un reparto que nadie quiera bloquear, esta arquitectura ofrece una configuración
que nadie \emph{pueda} bloquear sin coste físico inmediato. Es una respuesta de
ingeniería a una imposibilidad de teoría de juegos, y conviene enunciarla así en
lugar de dejar la impresión de que el reparto por precios resuelve la
estabilidad coalicional. No la resuelve; la traslada al contacto.

Queda con ello una cuarta lectura del mismo fenómeno, junto a
\cref{pr:caza-ciervo}, \cref{obs:deficit-capas} y \cref{obs:fuerte-ciervo}: la
formación de coaliciones es intrínsecamente inestable en el plano de los
incentivos, y todo lo que la sostiene en este trabajo —certificado, guarda,
acoplamiento— son mecanismos de compromiso, no de incentivo.
\end{observacion}

\subsection{Retribuir por aportación no es repartir con equidad}

Hay una tercera pregunta, distinta de las dos anteriores y que este trabajo
responde sin haberla planteado: una vez formada la coalición y producido el
excedente, ¿cómo se reparte? El mecanismo de precios da una respuesta
—cada AMR cobra el precio común por lo que aporta— pero esa no es la única
respuesta posible, y la teoría de la negociación ofrece otra
\cite{clempner2024optimization}: repartir el excedente sobre el punto de
desacuerdo, que aquí es el pago de no participar.

\begin{observacion}[Las dos respuestas difieren]
\label{obs:precio-vs-negociacion}
Con dos AMR de costes cuadráticos $k_1=1$ y $k_2=3$, demanda $W=10$ y pagos de
desacuerdo $2$ y $1$, el reparto eficiente del esfuerzo es $(7{,}5;\ 2{,}5)$ y el
precio común, igual al coste marginal, resulta $15$. El excedente total sobre el
desacuerdo es $72$, y se distribuye así:
\[
  \text{por precio: }(54{,}25;\ 17{,}75)
  \qquad\text{frente a}\qquad
  \text{por negociación: }(36;\ 36).
\]
No coinciden, y la brecha no es pequeña. Coinciden únicamente cuando costes y
pagos de desacuerdo son simétricos: con $k_1=k_2=2$ y desacuerdos iguales, ambos
criterios dan $(48{,}5;\ 48{,}5)$. Basta romper la simetría de costes,
manteniendo iguales los desacuerdos, para que vuelvan a separarse:
$(54{,}25;\ 16{,}75)$ frente a $(35{,}5;\ 35{,}5)$.
\end{observacion}

\begin{observacion}[Un tercer criterio, intermedio]
\label{obs:kalai}
Entre los dos anteriores existe un tercero, el de Kalai--Smorodinsky, que reparte
el excedente en proporción a la \emph{aspiración máxima} de cada parte —lo más que
cada una podría obtener si la otra se conformase con su punto de desacuerdo—
\cite{clempner2024optimization}. Sobre el mismo ejemplo: con aspiraciones iguales
coincide con el de Nash, $(36;36)$; si el AMR barato puede aspirar a más —por
ejemplo al $100\%$ frente al $45\%$ del caro— da $(49{,}66;\ 22{,}34)$, entre el
reparto competitivo $(54{,}25;\ 17{,}75)$ y el igualitario.

Los tres criterios responden a la misma pregunta con axiomas distintos, y la
elección no la resuelve la matemática. Lo que sí puede exigirse es que se declare
cuál se usa, que es lo que hace \cref{obs:criterio-reparto}.
\end{observacion}

\begin{observacion}[Shapley: un cuarto criterio, y premia otra cosa]
\label{obs:shapley}
Cuando el núcleo es vacío —que por \cref{pr:nucleo} es el caso normal— la
alternativa canónica es el valor de Shapley, que retribuye a cada miembro por su
contribución marginal promediada sobre todos los órdenes de formación
\cite{okada2026games}. Con AMR simétricos reparte a partes iguales y no aporta
nada; su interés está en el caso heterogéneo, que es el de este trabajo.

Con capacidades $(3;\ 2;\ 1;\ 1;\ 0{,}5;\ 0{,}5)$ y umbral $4$, el valor de
Shapley da $(0{,}450;\ 0{,}183;\ 0{,}133;\ 0{,}133;\ 0{,}050;\ 0{,}050)$, frente
al reparto proporcional a la capacidad
$(0{,}375;\ 0{,}250;\ 0{,}125;\ 0{,}125;\ 0{,}063;\ 0{,}063)$. El más fuerte cobra
\emph{más} que proporcional y el segundo \emph{menos}.

La razón es instructiva: Shapley premia ser \emph{pivote} —estar en muchas
coaliciones donde la propia presencia decide si se alcanza el umbral— y no tener
mucha fuerza. Un AMR mediano puede ser prescindible en casi toda combinación y
cobrar por debajo de su capacidad. Es un criterio con lectura física clara y
distinta de la del precio, que retribuye por esfuerzo aportado al margen.
\end{observacion}

\begin{observacion}[Qué criterio adopta este trabajo, y por qué]
\label{obs:criterio-reparto}
Conviene decirlo explícitamente, porque es una elección normativa y no un
teorema. El mecanismo de este trabajo es \emph{competitivo}: retribuye por
aportación al margen. El AMR barato aporta más y cobra más, y quien aporta poco
cobra poco aunque su presencia fuera necesaria para alcanzar el certificado.

La alternativa negociadora repartiría el excedente por igual sobre el
desacuerdo, con lo que un miembro imprescindible pero costoso no quedaría
penalizado por serlo. Tiene a su favor un argumento de estabilidad de coalición:
un miembro que se sienta mal retribuido tiene incentivo a abandonar, y por
\cref{pr:caza-ciervo} basta con que lo crea para que el equilibrio malo gane
terreno.

Se adopta el criterio competitivo porque es el que se deriva del dual del
programa convexo —los precios no se eligen, salen— y porque preserva la
estructura de potencial de \cref{th:potencial}, sobre la que descansa la
terminación. Un reparto negociado exigiría resolver un problema aparte en cada
formación y no está garantizado que conserve esa estructura. Se declara, pues,
como una decisión de diseño con un coste identificado, no como la única opción
razonable.
\end{observacion}

\subsection{El concepto que descarta el mal equilibrio}

\Cref{pr:caza-ciervo} exhibe dos equilibrios y \cref{obs:precio-no-basta}
muestra que ningún pago moderado elimina el malo. Existe un refinamiento del
concepto de equilibrio que sí lo descarta, y resulta ser el que corresponde a
este problema.

\begin{definicion}[Equilibrio fuerte]
\label{def:fuerte}
Un perfil es un \emph{equilibrio de Nash fuerte} si ninguna coalición de agentes
—de cualquier tamaño, no solo un agente aislado— puede desviarse conjuntamente
de modo que \emph{todos} sus miembros mejoren \cite{clempner2024optimization}.
\end{definicion}

\begin{observacion}[El mal equilibrio es de Nash pero no es fuerte]
\label{obs:fuerte-ciervo}
Con seis AMR, umbral $r^\star=4$, $B=10$, $c=1$ y $w=3$, la comprobación
exhaustiva sobre las $2^{6}$ configuraciones y todas las coaliciones da:
\begin{itemize}[leftmargin=1.9em,topsep=0.2em,itemsep=0.15em]
  \item «nadie se une»: es equilibrio de Nash —ningún AMR gana desviándose
  solo— pero \emph{no} es fuerte: una coalición de exactamente cuatro mejora a
  los cuatro, de $3$ a $9$;
  \item «todos se unen»: es equilibrio de Nash \emph{y} es fuerte.
\end{itemize}
El refinamiento separa exactamente los dos equilibrios que
\cref{pr:caza-ciervo} no distinguía.
\end{observacion}

\begin{observacion}[La guarda del autómata implementa la desviación de coalición]
\label{obs:guarda-fuerte}
La consecuencia es más concreta que un refinamiento conceptual. La desviación que
descalifica al mal equilibrio es la de una coalición de tamaño $r^\star$ que se
compromete \emph{conjuntamente}. Y eso es literalmente lo que hace la guarda de
\textsf{REC} a \textsf{FORM} de \cref{def:automata}: reunir exactamente los
miembros necesarios y comprometerlos a la vez.

Es decir, el certificado no solo hace creíble la cooperación
(\cref{obs:precio-no-basta}): \emph{es} el mecanismo que ejecuta la desviación
coalicional que el equilibrio fuerte exige. Un mecanismo que solo permitiese
revisiones individuales se quedaría atrapado en el mal equilibrio por
construcción, porque ninguna desviación unilateral es rentable allí.

Queda así una lectura común de tres hallazgos que habían aparecido por separado:
la multiplicidad de \cref{pr:caza-ciervo}, el déficit positivo de
\cref{obs:deficit-capas} y el fallo de fortaleza de \cref{obs:fuerte-ciervo} son
tres medidas del mismo fenómeno, y el compromiso conjunto es su único remedio
estructural.
\end{observacion}

\subsection{Subprovisión de la reserva}

Un AMR en reserva no transporta. Si una coalición sufre un fallo, un AMR
disponible cerca puede sustituir al miembro perdido y salvar la carga. El
beneficio de ese rescate lo recibe sobre todo \emph{un tercero} —la misión que
no queda varada—, mientras que el rescatador solo captura la parte que el
mecanismo le pague.

Sea $s\in[0,1]$ la fracción de un plantel de $N$ AMR que permanece en reserva,
$\Lambda$ el número de misiones activas, $\lambda_f$ la tasa de fallo por misión
y $p(s)$ la probabilidad de que un fallo encuentre rescate, creciente y cóncava
con $p(0)=0$. Sea $B$ el beneficio social de un rescate, $R$ el premio que cobra
el rescatador y $c$ el coste de oportunidad de estar en reserva.

\begin{supuesto}[Plantel de reserva separado]
\label{s:plantel}
$\Lambda$ no depende de $s$: la reserva se dota de un plantel distinto del que
ejecuta las misiones activas.
\end{supuesto}

La asimetría que genera la ineficiencia no está en el precio, sino en qué
cantidad ve cada parte. Los $Ns$ AMR en reserva son anónimos e intercambiables,
de modo que los $\Lambda\lambda_f p(s)$ rescates por unidad de tiempo se reparten
entre ellos por igual: cada uno ejecuta $\Lambda\lambda_f p(s)/(Ns)$, y percibe
por tanto la \emph{media}. El planificador, en cambio, evalúa el efecto de una
unidad más de reserva sobre el bienestar, que es la \emph{marginal} $p'(s)$.

\begin{teorema}[La brecha entre la media y la marginal]
\label{th:subprovision}
Bajo \cref{s:plantel}, sea $s^\star$ el nivel de equilibrio, definido por la
condición de entrada y salida
\[
  \frac{\Lambda\lambda_f\,p(s^\star)}{N\,s^\star}\,R\;=\;c,
\]
y sea $s^{\mathrm{soc}}$ el que maximiza $W(s)=\Lambda\lambda_f p(s)B-sNc$, es
decir $\Lambda\lambda_f p'(s^{\mathrm{soc}})B=Nc$. Sea
$\varepsilon(s)=s\,p'(s)/p(s)$ la elasticidad de $p$. Entonces:
\begin{enumerate}[label=(\roman*),leftmargin=2.2em,topsep=0.2em,itemsep=0.15em]
  \item $\varepsilon(s)\le1$ para todo $s>0$, con igualdad solo si $p$ es lineal:
  la media siempre domina a la marginal;
  \item $s^\star=s^{\mathrm{soc}}$ si y solo si $R=B\,\varepsilon(s^{\mathrm{soc}})$;
  \item $s^\star<s^{\mathrm{soc}}$ si y solo si $R<B\,\varepsilon(s^{\mathrm{soc}})$,
  y $s^\star>s^{\mathrm{soc}}$ si y solo si $R>B\,\varepsilon(s^{\mathrm{soc}})$.
\end{enumerate}
\end{teorema}

\begin{proof}
(i) Por concavidad $p'$ es decreciente, luego
$p(s)=\int_0^s p'(u)\,du\ge s\,p'(s)$, que es $\varepsilon(s)\le1$.
(ii) Imponiendo $s^\star=s^{\mathrm{soc}}$ en la condición de equilibrio y
sustituyendo $Nc=\Lambda\lambda_f p'(s^{\mathrm{soc}})B$ se obtiene
$R=B\,s^{\mathrm{soc}}p'(s^{\mathrm{soc}})/p(s^{\mathrm{soc}})=B\varepsilon(s^{\mathrm{soc}})$.
(iii) $p(s)/s$ es estrictamente decreciente por concavidad, luego $s^\star(R)$ es
estrictamente creciente en $R$; combinado con (ii) se sigue la equivalencia.
\end{proof}

\begin{observacion}[La reserva no emerge]
\label{obs:no-emerge}
Con $N=20$, $\Lambda=6$, $\lambda_f=0{,}05$, $c=1$, $B=40$, $R=12$ y
$p(s)=1-e^{-3s}$, resulta $s^\star=0$ y $s^{\mathrm{soc}}=0{,}196$, con pérdida
de bienestar $1{,}41$. Es decir: con esos parámetros \emph{ningún} AMR elige
reserva por su cuenta, aunque el óptimo social exija que casi uno de cada cinco
lo haga. Un banco de pruebas que asigne su reserva a mano no está tomando un
atajo de implementación: está compensando una ineficiencia del mecanismo.
\end{observacion}

\begin{observacion}[Pagar el beneficio completo sobreprovee]
\label{obs:sobreprovision}
De \cref{th:subprovision}(iii) se sigue un resultado contrario a la intuición:
entregar al rescatador el beneficio social íntegro, $R=B$, \emph{no} implementa
el óptimo, porque $B>B\varepsilon$ siempre que $p$ sea estrictamente cóncava.
Con los parámetros de \cref{obs:no-emerge} se obtiene
$\varepsilon(s^{\mathrm{soc}})=0{,}735$, de modo que $R=B=40$ conduce a
$s^\star=0{,}440$ frente a $s^{\mathrm{soc}}=0{,}196$: un $124\%$ de exceso de
reserva. La razón es que un rescatador anónimo cobra por cada rescate que
ejecuta —la media— mientras que su aportación al bienestar es solo la marginal.
Corregir la externalidad no consiste en trasladar todo el beneficio, sino
exactamente la fracción $\varepsilon$.
\end{observacion}

\subsection{El cuarto precio}

La corrección es la internalización habitual de una externalidad positiva, y
encaja en la arquitectura sin añadir capas: es \emph{otro precio}.

\begin{teorema}[Implementación por prima de cobertura]
\label{th:prima}
Fíjese el premio del rescatador en
$\widetilde R_\ell=B_\ell\,\varepsilon(s^{\mathrm{soc}})$ y hágase pagar a cada
misión activa una prima por unidad de tiempo
\[
  q_\ell\ =\ \lambda_f\,p(s^{\mathrm{soc}})\,\widetilde R_\ell
  \ =\ \lambda_f\,p(s^{\mathrm{soc}})\,B_\ell\,\varepsilon(s^{\mathrm{soc}}).
\]
Entonces: (i) el equilibrio verifica $s^\star=s^{\mathrm{soc}}$; (ii) el
mecanismo está equilibrado en esperanza, es decir, las primas recaudadas igualan
los pagos esperados; (iii) la prima no excede el valor esperado de la cobertura
para la misión, $\lambda_f p(s^{\mathrm{soc}})B_\ell$, luego la participación es
individualmente racional.
\end{teorema}

\begin{proof}
(i) es \cref{th:subprovision}(ii). (ii) La recaudación por unidad de tiempo es
$\Lambda\lambda_f p(s^{\mathrm{soc}})\widetilde R$ y el pago esperado es el
mismo: hay $\Lambda\lambda_f$ fallos por unidad de tiempo, cada uno rescatado
con probabilidad $p(s^{\mathrm{soc}})$ y retribuido con $\widetilde R$. (iii) es
$\varepsilon\le1$ de \cref{th:subprovision}(i).
\end{proof}

\begin{observacion}[Con aversión al riesgo el margen es mucho mayor]
\label{obs:aversion-riesgo}
El apartado (iii) demuestra racionalidad individual para participantes
\emph{neutrales al riesgo}: la prima no excede el valor esperado de la cobertura.
Una flota real no es neutral —una misión varada tiene costes de reputación y de
compromiso que crecen más que linealmente con la pérdida—, y bajo aversión al
riesgo el resultado se refuerza en lugar de debilitarse
\cite{clempner2024optimization}.

Con utilidad exponencial y los parámetros de \cref{obs:no-emerge}, la prima
cobrada es $13{,}065$ y el valor esperado de la cobertura $17{,}776$, de modo que
la holgura ya es de $+4{,}71$ con neutralidad. Al aumentar la aversión, la
disposición máxima a pagar sube a $19{,}71$, $23{,}16$, $26{,}92$ y $30{,}71$
para coeficientes de $0{,}01$ a $0{,}12$, con holguras de $+6{,}6$ hasta
$+17{,}6$.

Hay además una consecuencia de diseño que no se sigue de la versión neutral: si
quien opera la flota es averso al riesgo, su nivel \emph{óptimo} de reserva
excede el $s^{\mathrm{soc}}$ de \cref{th:subprovision}, porque ese cálculo
maximiza el bienestar esperado y no incorpora prima de riesgo alguna. El nivel de
\cref{th:subprovision} debe leerse, por tanto, como una cota inferior del que un
operador prudente elegiría.
\end{observacion}

\begin{observacion}[Cuánto cambia el número]
\label{obs:prima-num}
Con los parámetros de \cref{obs:no-emerge} la prima correcta es
$q=\lambda_f p(s^{\mathrm{soc}})B\varepsilon=0{,}653$ por unidad de tiempo, frente
a $0{,}889$ que resultaría de retribuir el rescate con $B$ íntegro. La diferencia
no es un ajuste fino: es la que separa el óptimo de un exceso de reserva del
$124\%$, con su coste de oportunidad correspondiente.
\end{observacion}

\begin{observacion}[La prima es el pago esperado, no la pérdida retenida]
\label{obs:prima}
Conviene no confundir tres cantidades. La \emph{pérdida esperada} de una misión
es $\lambda_f(1-p(s))B_\ell$: lo que se pierde cuando el rescate \emph{no} llega.
El \emph{valor esperado de la cobertura} es $\lambda_f p(s)B_\ell$: lo que la
misión gana cuando el rescate \emph{sí} llega. Y la \emph{prima} que financia el
mecanismo es $\lambda_f p(s)B_\ell\varepsilon$, estrictamente menor que la
segunda por \cref{th:subprovision}(i). Las dos primeras coinciden solo si
$p=1/2$; confundir cualquiera de ellas con la tercera desequilibra el fondo o
distorsiona el nivel de reserva. Que la prima quede por debajo del valor de la
cobertura es lo que hace el mecanismo aceptable para quien paga.
\end{observacion}

Con esto el mecanismo queda con cuatro familias de precios, todas obtenidas del
mismo modo —restricción o externalidad, programa convexo, dual devuelto como
pago— y todas difundidas por el mismo consenso:

\begin{table}[!ht]
\centering\small
\caption{Las cuatro familias de precios y qué internaliza cada una.}
\label{tab:precios}
\begin{tabularx}{\textwidth}{@{}llX@{}}
\toprule
Precio & Recurso & Origen \\
\midrule
$p_k$ & capacidad de \emph{wrench} & dual del reparto \eqref{eq:qp} \\
$\tau^c_e$ & espacio del corredor & externalidad de congestión por clase (\cref{th:multiclase}) \\
$\lambda_{ij}$ & seguridad & multiplicador de la barrera (\cref{th:precio-seguridad}) \\
$q_\ell$ & riesgo & externalidad de rescate (\cref{th:prima}) \\
\bottomrule
\end{tabularx}
\end{table}

\begin{observacion}[Los cuatro precios son una escalarización, y eso tiene un límite]
\label{obs:escalarizacion}
Los cuatro precios convierten un problema con varios objetivos —caudal, energía,
espacio, riesgo— en uno solo, sumando cada término con su peso. Esa operación
tiene una limitación conocida y conviene declararla: la escalarización lineal
alcanza únicamente los puntos del frente de Pareto que están sobre la envolvente
convexa de la región factible. Los tramos no convexos del frente son
inalcanzables \emph{para cualquier elección de precios}, no solo para una mala
elección \cite{clempner2024optimization}.

La comprobación es contundente. Sobre un frente con un tramo cóncavo,
discretizado en $4001$ puntos, un barrido de cuatrocientos juegos de pesos
alcanza $284$ puntos, el $7{,}1\%$. El resto del frente es Pareto-óptimo y
ningún precio lo selecciona.

Esto no invalida el mecanismo, pero acota lo que puede afirmarse de él. Cuando
se dice que los precios internalizan las externalidades y conducen a un óptimo,
se está diciendo que conducen al óptimo \emph{de la combinación ponderada
elegida}. Si la frontera de compromisos real tiene tramos no convexos —y una
región factible definida por restricciones combinatorias como el catálogo de
coaliciones los tiene con facilidad—, existen puntos de operación eficientes que
la arquitectura no puede alcanzar por ningún ajuste de precios. Alcanzarlos
exigiría un método distinto, de los que trabajan sobre el frente en lugar de
sobre una suma ponderada.
\end{observacion}

\begin{observacion}[Quién mueve primero]
\label{obs:stackelberg}
Conviene además nombrar la estructura temporal del mecanismo, que hasta aquí ha
quedado implícita. Los peajes y la prima no se negocian con los AMR: se fijan y
los AMR responden a ellos. Esa es una estructura de \emph{líder y seguidores}
—de Stackelberg— y no un juego simultáneo entre iguales
\cite{clempner2024optimization}. El diseñador anticipa la mejor respuesta de la
población y elige el precio que la conduce al óptimo social, que es exactamente
lo que hacen \cref{th:multiclase} y \cref{th:prima}.

La precisión importa porque el equilibrio de Stackelberg y el de Nash del juego
simultáneo no coinciden en general, y porque la capacidad del líder para
comprometerse con un precio anunciado es una hipótesis, no un dato: si los AMR
esperasen que el precio se renegocie tras ver su respuesta, el mecanismo pierde
su fuerza. En esta arquitectura el compromiso lo proporciona el hecho de que el
precio sea el dual de un programa público y verificable, no una decisión
discrecional.
\end{observacion}

\begin{observacion}[Atribución]
\label{obs:standby-atrib}
La subprovisión de un bien público y su corrección pigouviana son teoría
económica estándar, y el caso binario con beneficio a terceros es el
\emph{dilema del voluntario}. Lo que corresponde a este trabajo es el dominio y
la composición: la reserva de una flota como bien público \emph{espacial}
—$p(s)$ depende de la geometría, no solo del número—, con la prima difundida
por el mismo mecanismo de consenso que los otros tres precios y verificable
localmente, porque el rescate es un evento físico observable por el vecindario
que lo presencia. No se introduce ningún mecanismo nuevo: se añade un precio a
un mercado que ya tenía tres.
\end{observacion}

\subsection{Conectividad factorizada}

\Cref{th:topologia-garantizada} impone un árbol generador sobre \emph{toda} la
flota y lo mantiene indefinidamente. Es más de lo necesario, y su coste
—movilidad restringida— es real. La estructura de misiones permite factorizar.

\begin{proposicion}[Conectividad por fases]
\label{pr:factorizada}
Basta con: (i) conectividad \emph{intra}-coalición permanente, que la propia
formación garantiza por geometría cuando su diámetro no excede
$R^{\mathrm{com}}$; y (ii) conectividad \emph{inter}-coalición únicamente
durante la fase de reclutamiento, mientras el residuo del reparto está por
encima del umbral de compromiso. No se requiere conectividad global permanente.
\end{proposicion}

\begin{proof}
Los precios de \emph{wrench} y de seguridad de una coalición solo involucran a
sus miembros, luego su consenso vive en el subgrafo intra-coalición, conexo por
(i). El consenso de reclutamiento y de precios de carga solo debe converger
mientras la decisión sigue abierta; una vez comprometido el perfil, su valor no
se consulta. Aplicando \cref{th:conmutada} a la ventana de reclutamiento se
obtiene el acuerdo requerido.
\end{proof}

\begin{observacion}[La condición de carrera]
\label{obs:carrera}
\Cref{pr:factorizada} traslada la exigencia a una comparación de tiempos: el
reclutamiento debe converger antes de que las coaliciones se separen,
\[
  T_{\mathrm{conv}}\ <\ T_{\mathrm{desconexión}} .
\]
El primero se acota con la tasa de \cref{th:mapa}; el segundo, con la cinemática
de separación y el radio $R^{\mathrm{com}}$. No se afirma que la desigualdad se
cumpla siempre: es una condición verificable por instancia, y su violación
significa que el compromiso debe adelantarse o la separación retrasarse. El
ahorro respecto a \cref{th:topologia-garantizada} es que la restricción de
movilidad se paga solo durante el reclutamiento y solo entre coaliciones.
\end{observacion}

\subsection{Dónde colocar la reserva}
\label{sec:cobertura}

\Cref{th:subprovision} da \emph{cuánta} reserva sostener y \cref{th:prima} el
mecanismo que la sostiene, pero no dónde situarla. Esa pregunta —el
emplazamiento de vehículos ociosos, en el vocabulario de
\cref{tab:posicionamiento}— quedó declarada fuera. Se cierra aquí, porque
resulta ser un problema de cobertura y porque la heterogeneidad que atraviesa
todo el trabajo también decide su forma.

\begin{definicion}[Alcance de rescate y cobertura]
\label{def:cobertura}
Sea $\varphi:\Omega\to\Rb_{\ge0}$ la densidad espacial de fallos, es decir la
tasa con que se producen averías por unidad de área. Sea
$\Rcal_i=\{q\in\Omega:\lVert q-q_i\rVert\le\varrho_i\}$ el \emph{alcance de
rescate} del AMR $i$ situado en $q_i$, con
$\varrho_i=v^{\max}_i\,T^{\mathrm{res}}$ el radio que puede recorrer dentro del
plazo admisible. La cobertura de rescate es
\[
  \Hcal(q_1,\dots,q_n)\;=\;\int_\Omega\ \max_{i}\ \mathbf 1_{\Rcal_i}(q)\ \varphi(q)\,dq .
\]
\end{definicion}

\begin{observacion}[La cobertura es lo que hace de $p(s)$ una función]
\label{obs:cobertura-p}
$\Hcal$ normalizada por $\int_\Omega\varphi$ es exactamente la probabilidad de
que un fallo encuentre rescate a tiempo, es decir la función $p$ de
\cref{th:subprovision}, que hasta aquí se había postulado creciente y cóncava
sin decir de qué depende. Depende del emplazamiento: dos flotas con la misma
fracción $s$ en reserva pueden tener coberturas muy distintas. Nivel y
emplazamiento no son decisiones separables, y \cref{def:cobertura} es el puente
entre ambas.
\end{observacion}

\begin{teorema}[Con AMR heterogéneos la partición correcta no es la de Voronoi]
\label{th:potencia}
Asígnese a cada AMR la región de la que se responsabiliza. Si los alcances
$\varrho_i$ difieren, la partición de Voronoi —por proximidad— asigna puntos a
AMR que no llegan a tiempo mientras otro sí llegaría. La partición que respeta
los alcances es el \emph{diagrama de potencia}
\[
  \Pcal_i=\bigl\{q\in\Omega:\ \lVert q-q_i\rVert^2-\varrho_i^{\,2}
  \ \le\ \lVert q-q_j\rVert^2-\varrho_j^{\,2},\ \ \forall j\neq i\bigr\},
\]
que es una teselación completa de $\Omega$ en celdas convexas y coincide con la
de Voronoi si y solo si todos los alcances son iguales.
\end{teorema}

\begin{proof}
La convexidad y la teselación se siguen de que las fronteras son hiperplanos:
la condición define semiplanos, pues los términos cuadráticos $\lVert q\rVert^2$
se cancelan. La coincidencia con Voronoi cuando $\varrho_i\equiv\varrho$ es
inmediata. Para la inadecuación basta un contraejemplo, y se da en
\cref{obs:voronoi-falla}.
\end{proof}

\begin{observacion}[El contraejemplo]
\label{obs:voronoi-falla}
Con $q_1=(0,0)$, $\varrho_1=3$ y $q_2=(4,0)$, $\varrho_2=1$, el punto $q=(2{,}5,0)$
dista $2{,}5$ del primero y $1{,}5$ del segundo: Voronoi lo asigna al segundo,
que \emph{no} llega —$1{,}5>1$—, mientras el primero sí llega —$2{,}5\le3$—. El
diagrama de potencia lo asigna al primero, porque $2{,}5^2-3^2=-2{,}75$ es menor
que $1{,}5^2-1^2=1{,}25$. Lo mismo ocurre en $q=(2,0)$.

Que la partición correcta dependa de los alcances y no de las distancias es
\cref{th:noequiv} reapareciendo en la capa espacial: los AMR no son
intercambiables, y cualquier construcción que los trate como puntos idénticos
—aquí, una partición por proximidad— reparte mal.
\end{observacion}

\begin{proposicion}[Reubicación distribuida con mejora garantizada]
\label{pr:ascenso}
Con la ley $\dot q_i=\alpha_i\,\partial\Hcal/\partial q_i$, $\alpha_i>0$, se
tiene $\dot{\Hcal}=\sum_i\alpha_i\lVert\partial\Hcal/\partial q_i\rVert^2\ge0$:
la cobertura crece de forma monótona. Además $\partial\Hcal/\partial q_i$ se
evalúa como una integral sobre la frontera de $\Rcal_i$ dentro de la celda
$\Pcal_i$, luego depende solo de los vecinos en el grafo dual y la ley es
distribuida.
\end{proposicion}

\begin{observacion}[La monotonía no sobrevive al muestreo sin cota de paso]
\label{obs:ascenso-muestreo}
\Cref{pr:ascenso} es un enunciado de tiempo continuo, y conviene no leerlo como
si valiera para la implementación. Sobre una cobertura suavizada con tres AMR y
ochenta pasos, un paso $\alpha=0{,}10$ da monotonía estricta —cero descensos—;
$\alpha=0{,}30$ produce cuatro descensos; y $\alpha\ge0{,}80$, treinta y nueve,
aunque la cobertura final siga siendo mayor que la inicial. Es la misma
situación que \cref{le:muestras}: una garantía de tiempo continuo exige una cota
de paso explícita para trasladarse a un sistema muestreado, y omitirla convierte
un teorema en una expectativa.
\end{observacion}

\begin{observacion}[Incertidumbre de posición: se encoge el alcance]
\label{obs:cobertura-incertidumbre}
Si la pose del AMR se conoce solo dentro de una región de incertidumbre de radio
$r_i$, la cobertura \emph{garantizada} se obtiene sustituyendo $\varrho_i$ por
$\varrho_i-r_i$. Es la misma operación que \cref{co:tensado} realiza sobre las
barreras, con $r_i$ derivado de la misma covarianza del filtro. Aparece así por
tercera vez el mismo patrón: una garantía se conserva bajo incertidumbre
encogiendo el conjunto sobre el que se afirma, y el encogimiento se paga en
prestaciones —aquí, en área cubierta— de forma continua y no catastrófica.
\end{observacion}
